\section{الفصل~\ref{sec:neurontypes}. أنواع العصبونات}

تستند أرقام الفصل إلى عدّ مباشر للخلايا في دماغ الإنسان حيثما وُجد هذا العدّ، وتستند قاعدة ”عصبون واحد، ناقل واحد“ إلى الأطالس الخلوية وإلى تجارب وجدت استثناءاتها أيضًا، ويستند دور \nt{DOPA} إلى تسجيلات النبضات، أما أدوار قنوات البثّ الأخرى فتستند إلى فرضيات.

{\footnotesize
\setlength{\tabcolsep}{3pt}\renewcommand{\arraystretch}{1.15}
\begin{longtable}{>{\raggedright}p{0.5cm}>{\raggedright}p{4.0cm}>{\raggedright}p{5.6cm}>{\raggedright}p{2.3cm}>{\raggedright\arraybackslash}p{1.7cm}}
\toprule
م & القضية & التجربة وما قيس & المصدر & الحالة \\
\midrule\endfirsthead
\toprule
م & القضية & التجربة وما قيس & المصدر & الحالة \\
\midrule\endhead
1 & في دماغ الإنسان نحو $8.6\cdot10^{10}$ عصبون، وكل عصبونات \nt{GLUT} تقريبًا عصبونات صغيرة في المخيخ (الجدول~\ref{tab:types}) & المجزِّئ المتماثل الخواص: يُذاب النسيج إلى معلّق متجانس من النوى، وتُعدّ النوى الحاملة للواسم العصبوني \foreignlanguage{english}{NeuN}. عند الرجال البالغين $86.1\pm8.1$ مليار عصبون؛ منها $19\,\%$ فقط في قشرة نصفي الكرة المخيين، والكتلة الرئيسية في المخيخ & \cite{Azevedo2009} & مُقاس \\
2 & لكل عصبون تقريبًا ناقل عصبي رئيسي واحد (القضية~\ref{prop:dale})، لكن ثمة استثناءات & أطلس ترنسكريبتومي لدماغ الفأر: نحو $4\cdot10^6$ خلية، و$5322$ عنقودًا؛ نُسب إلى كل عنقود ناقله بحسب جينات التخليق والتعبئة. والاستثناءات مقيسة مباشرة: عصبونات \nt{DOPA} في الشرائح تقذف الغابا أيضًا عبر الناقل الحويصلي نفسه وتثبط عصبونات المخطط؛ وفي جزء من محاور \nt{DOPA} يخرج الغلوتامات والدوبامين من نهايات مختلفة للسلك نفسه (المجهر الإلكتروني وعلم البصريات الوراثي) & \cite{Strata1999,Yao2023,Tritsch2012,Zhang2015} & مُقاس \\
3 & العمود كله في مصفوفة الأوزان بعلامة واحدة~\eqref{eq:dalesign} & استنتاج في الفصل من القضية~\ref{prop:dale} ومن أن مستقبلات الغلوتامات كلها تقريبًا مثيرة، ومستقبلات الغابا مثبطة. لا يوجد اختبار مباشر لقاعدة ”كل مستقبِلي العصبون الواحد يتلقون وزنًا بعلامة واحدة“ بوصفها قاعدة عامة؛ ومن الأمثلة المضادة القذف المشترك للغابا من عصبونات \nt{DOPA} (السطر~2) والمستقبِلون الحاملون لمستقبلات مختلفة العلامة (السطر~11) & استنتاج في الفصل & نظرية \\
4 & من ثلاثة أرباع إلى أربعة أخماس عصبونات القشرة من \nt{GLUT}، ومن الخُمس إلى الربع من \nt{GABA} & صبغ مناعي للغابا في أعمدة عرضها $50\um$ تخترق سماكة القشرة كلها في $10$ مناطق من قشرة القرد: في $8$ مناطق تشكل عصبونات الغابا نحو $25\,\%$ من كل العصبونات، وفي القشرة البصرية الأولية أقل من $20\,\%$ & \cite{Hendry1987} & مُقاس \\
5 & لعصبون \nt{GLUT} نحو $10^4$ مخرج & علم القياس المجسم عند التشريح: في القشرة الحديثة للرجال الشباب $1.64\cdot10^{14}$ مشبك (المجهر الإلكتروني، طريقة الشريحتين) ونحو $2.3\cdot10^{10}$ عصبون. والنسبة تعطي $\sim7\cdot10^3$ مشبك لكل عصبون؛ وفي المتوسط يساوي عدد المداخل عدد المخارج & \cite{Tang2001synapses,Pakkenberg1997} & مُقاس \\
6 & يتفرع مخرج عصبون \nt{DOPA} إلى $10^5\text{--}10^6$ نهاية؛ وهذا تقدير من مدى تغطية الهدف لا عدّ & وسم فيروسي لغشاء عصبونات \nt{DOPA} منفردة لدى الجرذ وإعادة بناء كاملة للمحور: العصبون الواحد يغطي من $0.45$ إلى $5.7\,\%$ (في المتوسط $2.7\,\%$) من حجم المخطط. قيست التغطية لا عدد النهايات: فهذا مستنتج من التغطية من حيث رتبة المقدار، ولا عدّ مباشرًا للنهايات. وأعداد النهايات لـ\nt{SERO} و\nt{NORA} تقديرات خشنة & \cite{Matsuda2009} & مُقاس \\
7 & عصبونات البثّ قليلة: \nt{DOPA} $\sim5\cdot10^5$ (المجموعة الرئيسية من مجموعتين، حدّ أدنى)، \nt{SERO} $\sim3\cdot10^5$ (مجموع عدّين جزئيين)، \nt{HIST} $\sim6\cdot10^4$ (الجدول~\ref{tab:types}، الشكل~\ref{fig:typepie}) & عدّ مجسم عند الإنسان: $550\,000$ عصبون مصطبغ في المادة السوداء (من دون السقيفة البطنية)؛ و$235\,000$ عصبون في نواة الرفاء الظهرية (كل عصبونات النواة)، و$125\,000$ عصبون سيروتونيني آخر في سقيفة الجسر؛ ونحو $64\,000$ عصبون هستاميني في الوطاء. لا عدّ مباشرًا لـ\nt{NORA} و\nt{GLYC} و\nt{ACET} و\nt{ADRE}: فهي في الجدول تقديرات خشنة لرتبة المقدار & \cite{Pakkenberg1991,Baker1990,Baker1991,Airaksinen1991} & مُقاس \\
8 & يقفز الغلوتامات في الشق إلى نحو ميلي مول ويهبط في $\sim1\ms$ & بإزاحة حاصر تنافسي سريع الانفصال عن مستقبلات \foreignlanguage{english}{NMDA} أثناء النقل المشبكي (مزرعة عصبونات الحُصين): الذروة $1.1\mM$، وثابت الهبوط $1.2\ms$ & \cite{Clements1992} & مُقاس \\
9 & في القشرة العاملة يكاد التياران المثير والمثبط يتوازنان، ويبقى العصبون قرب العتبة & النظرية: الشبكة ذات الوصلات القوية تبلغ النمط المتوازن من تلقاء نفسها. التجربة: في القشرة الجبهية الأمامية لابن مقرض في الكائن الحي، أثناء حالات ”الصعود“، يبقى جهد انعكاس التيار المشبكي قرب $-37\mV$ مئات من أجزاء الألف من الثانية، مع أن الموصلية تتغير بـ$\sim21\,\text{nS}$: الإثارة والتثبيط يرتفعان وينخفضان معًا & \cite{vanVreeswijk1996,Haider2006} & مُقاس \\
10 & عصبونات \nt{DOPA} تبلّغ عن خطأ التنبؤ بالمكافأة~\eqref{eq:rpe} (الشكل~\ref{fig:rpe}) & نبضات عصبونات \nt{DOPA} لدى القرد: رشقة للعصير غير المتوقع؛ وبعد التعلم رشقة للإشارة ولا استجابة للعصير المتنبأ به؛ وهبوط في التردد حين لا يُعطى العصير الموعود. وفي تجربة كمية يتناسب التردد مع الفرق بين المكافأة الحالية والمتوسط الموزون للمكافآت السابقة، لكن للأخطاء الموجبة فقط. أما المعادلة~\eqref{eq:rpe} نفسها فهي خوارزمية الفروق الزمنية & \cite{Schultz1997,Bayer2005,SuttonBarto2018} & مُقاس \\
11 & العلامة والسرعة يحددهما المستقبِل: المستقبلات مؤينة التوجه في أجزاء من جزء الألف من الثانية، وأيضية التوجه أبطأ بكثير (القضية~\ref{prop:receptor}) & نبضات سريعة من الغلوتامات على قطع من غشاء عصبونات الحُصين: التيار عبر المستقبلات مؤينة التوجه يرتفع (من $20$ إلى $80\,\%$) في $0.2\text{--}0.6\ms$ ويهبط في $2\text{--}3\ms$. والجهد المثبط البطيء للعصبونات الهرمية يزيله حاصر انتقائي للمستقبل أيضي التوجه للغابا. عائلتا مستقبلات الدوبامين متعاكستا التأثير؛ وفي المخطط تحتاج العصبونات ذات المستقبل \foreignlanguage{english}{D1} إلى الدوبامين لتقوية المشابك، وذات \foreignlanguage{english}{D2} لإضعافها & \cite{Colquhoun1992,DutarNicoll1988,Beaulieu2011,Shen2008} & مُقاس \\
12 & قناة البثّ ليست سلكًا واحدًا، بل حزمة من خطوط قليلة (القسم~\ref{sec:buses}) & وسم فيروسي-وراثي لعصبونات السيروتونين في نواة الرفاء الظهرية للفأر: العصبونات التي ترسل مخرجها إلى اللوزة وتلك التي ترسله إلى القشرة الجبهية تقع في مواضع مختلفة من النواة، وتتفرع إلى مناطق متتامة، وتستجيب استجابتين متعاكستين لمنبه مزعج. مراجعة: مجموعات فرعية من عصبونات الموضع الأزرق (\nt{NORA}) ترمّز عمليات مختلفة & \cite{Ren2018,Poe2020} & مُقاس \\
13 & \nt{NORA} يحدد الكسب $g$ & فرضية الكسب التكيفي؛ ونموذج شبكة ترفع فيه الكاتيكولامينات ميل منحنى النقل. لا قياس مباشرًا لـ”ازدياد $g$ مع النورأدرينالين“ في الشبكة كلها؛ وقد قيس أن قطر الحدقة يتبع نبضات عصبونات الموضع الأزرق لدى القرد & \cite{AstonJones2005,ServanSchreiber1990,Joshi2016} & فرضية \\
14 & \nt{ACET} يبدّل ”كتابة/قراءة“ ويحدد سرعة التعلم & فرضية. سندها: في شرائح القشرة الشمية للجرذ تضعف ناهضات الأسيتيل كولين ($100\,\mu\text{M}$) المشابك الداخلية ”المستحضِرة“ بأكثر من $60\,\%$، والمشابك الداخلة بأقل من $15\,\%$ & \cite{Hasselmo2006,HasselmoBower1992} & فرضية \\
15 & \nt{SERO} يحدد معامل الخصم $\gamma$؛ والعصبونات السيروتونينية تعمل بانتظام، ببضع نبضات في الثانية، وتكاد تصمت في إحدى مراحل النوم & فرضية ”السيروتونين $=\gamma$“. أقوى حجة: التنشيط البصري الوراثي لعصبونات السيروتونين في نواة الرفاء الظهرية للفأر يقلل عدد حالات التخلي المبكر عن انتظار مكافأة مؤجلة. وتردد النبضات والصمت في نوم حركة العين السريعة مقيسان (مراجعة للتسجيلات) & \cite{Doya2002,Miyazaki2014,JacobsAzmitia1992} & فرضية \\
\bottomrule
\end{longtable}}
